\section{Construction of $n$-Dependent Ladder Operators}

Although the infinite square well does not admit standard Heisenberg ladder operators with constant energy spacing, it is nevertheless possible to construct \emph{generalized} or \emph{nonlinear} creation and annihilation operators that act on the eigenbasis of the Hamiltonian. These operators connect consecutive energy levels with coefficients that depend explicitly on the quantum number $n$, reflecting the quadratic nature of the spectrum \cite{Hussin2011, Jalali2013, Fernandez1999}.

\subsection*{1. Definition on the Spectral Basis}

Let $\{|n\rangle\}_{n\ge1}$ denote the orthonormal eigenbasis of the single-particle Hamiltonian,
\[
\hat{H}|n\rangle = E_n |n\rangle, \qquad E_n = \frac{\pi^2\hbar^2}{2mL^2}n^2.
\]
We define the generalized ladder operators by their matrix elements:
\begin{equation}
    \hat{A}^\dagger = \sum_{n=1}^{\infty} c_n\,|n+1\rangle\langle n|, \qquad
    \hat{A} = \sum_{n=1}^{\infty} c_n^*\,|n\rangle\langle n+1|.
\end{equation}
Their action on the eigenbasis is
\begin{equation}
    \hat{A}^\dagger|n\rangle = c_n|n+1\rangle, \qquad
    \hat{A}|n\rangle = c_{n-1}^*|n-1\rangle,
\end{equation}
where the coefficients $c_n$ are chosen according to the spectral structure of the system.

\subsection*{2. Commutation Relation with the Hamiltonian}

Applying $\hat{H}$ to $\hat{A}^\dagger$ gives
\begin{equation}
    [\hat{H}, \hat{A}^\dagger] = \sum_{n=1}^{\infty} (E_{n+1} - E_n)\,c_n\,|n+1\rangle\langle n|.
\end{equation}
In contrast with the harmonic oscillator, where $(E_{n+1}-E_n) = \hbar\omega$ is constant, here the difference depends on $n$:
\begin{equation}
    E_{n+1} - E_n = \frac{\pi^2\hbar^2}{2mL^2}(2n+1).
\end{equation}
Hence, no constant $\lambda$ satisfies $[\hat{H},\hat{A}^\dagger] = \lambda \hat{A}^\dagger$.  
Nevertheless, one can express the commutator in the \emph{functional form}
\begin{equation}
    [\hat{H},\hat{A}^\dagger] = F(\hat{N})\,\hat{A}^\dagger,
\end{equation}
where $\hat{N}|n\rangle = n|n\rangle$ and $F(\hat{N})$ is the nonlinear step operator
\[
F(\hat{N}) = \frac{\pi^2\hbar^2}{2mL^2}\,(2\hat{N}+1).
\]
This defines a \textbf{deformed Heisenberg algebra}, also referred to as an \emph{$n$-dependent ladder algebra}, characterized by variable energy spacing.

\subsection*{3. Choice of Coefficients}

A convenient and physically motivated choice for the coefficients is
\begin{equation}
    c_n = \sqrt{2n+1},
\end{equation}
which mirrors the scaling of the energy differences $E_{n+1}-E_n$.  
This yields a consistent normalization and ensures that the operators $\hat{A}$ and $\hat{A}^\dagger$ correctly reproduce the spectral hierarchy:
\[
\hat{A}^\dagger |n\rangle \propto |n+1\rangle, \qquad \hat{A} |n\rangle \propto |n-1\rangle.
\]
The resulting algebra obeys
\begin{equation}
    [\hat{A}, \hat{A}^\dagger] = G(\hat{N}),
\end{equation}
where $G(\hat{N}) = F(\hat{N}) - F(\hat{N}-1)$ is a nonlinear function of the number operator.

\subsection*{4. Two-Particle Operators}

For a system of two indistinguishable particles, the total Hamiltonian is additive:
\[
\hat{H} = \hat{H}_1 + \hat{H}_2.
\]
One then defines particle-local operators
\[
\hat{A}_1^\dagger = \hat{A}^\dagger \otimes I, \qquad
\hat{A}_2^\dagger = I \otimes \hat{A}^\dagger,
\]
and constructs collective (anti)symmetrized combinations that preserve exchange symmetry:
\[
\hat{A}^\dagger_{\text{asym}} = \frac{1}{\sqrt{2}}\left(\hat{A}_1^\dagger - \hat{A}_2^\dagger\right), \qquad
\hat{A}^\dagger_{\text{sym}} = \frac{1}{\sqrt{2}}\left(\hat{A}_1^\dagger + \hat{A}_2^\dagger\right).
\]
For fermions, the antisymmetric operator $\hat{A}^\dagger_{\text{asym}}$ acts as
\[
\hat{A}^\dagger_{\text{asym}}|\Psi_F\rangle \propto 
\frac{1}{\sqrt{2}}\Big(|n_1+1, n_2\rangle - |n_2, n_1+1\rangle\Big),
\]
which preserves the antisymmetry required by the exchange statistics.

\subsection*{5. Discussion and Conclusion}

The generalized ladder operators $\hat{A}$ and $\hat{A}^\dagger$ constructed in this manner \emph{do exist} and act consistently on the eigenbasis of the infinite square well.  
However, their algebraic structure is inherently nonlinear:
\[
[\hat{H},\hat{A}^\dagger] = F(\hat{N})\hat{A}^\dagger, \qquad [\hat{A},\hat{A}^\dagger] = G(\hat{N}),
\]
where $F$ and $G$ depend on $\hat{N}$.  
Therefore, the infinite well does not support a standard Heisenberg algebra with constant spacing, but it does admit an \textbf{$n$-dependent (deformed) algebraic realization}.  
In this sense, ladder operators can indeed be constructed for the system, provided one allows their action and commutation relations to depend explicitly on the quantum number $n$.  

This construction aligns with the generalized Heisenberg algebras developed in \cite{Hussin2011, Jalali2013, Fernandez1999}, where systems with nonlinear spectra---such as the infinite well, Pöschl–Teller, and Morse potentials---are successfully described using $n$-dependent ladder operators that obey deformed commutation relations.


